\section{Rank loss on the boundary and the ramification divisor}
\label{sec:rankloss}

The preceding lemmas apply to a smooth model of a full period image.
At a neat level the variation is unipotent at the boundary. Returning
to the original level also introduces ramification divisors. We now
prove the strict numerical rank drop for both kinds of divisor.

We specify the finite-group data because the group acting on the
variation can be larger than the effective group acting on its base.
Let $T$ be a smooth connected projective complex variety of positive
dimension, let $D_T$ be a reduced simple-normal-crossing divisor, and
put $U=T\setminus D_T$. Let $\mathbb V$ be a rationally polarized
integral pure variation on $U$, minimal for its entire highest
Hodge line, with unipotent boundary monodromy. Write $J$ for the
extended highest line. Let $D_N$ be the union of the components
of $D_T$ with nonzero monodromy logarithm, and put
$U_N=T\setminus D_N$. The variation extends purely over $U_N$ by
Theorem~\ref{altord:h:inputs}.

Suppose a finite group $H$ acts on $T$, preserves $D_T$, and acts on
the variation by flat Hodge isomorphisms preserving its integral
lattice and rational polarization and covering that action. Set
\[
                 G_f=H/\ker\bigl(H\longrightarrow\Aut(T)\bigr).
\]
Thus $G_f$ acts effectively on $T$. We retain $H$ for its action on
the variation, since its kernel on $T$ can act nontrivially on fibres.
Uniqueness of the pure extension preserves this action over $U_N$.

Fix a connected period domain $\mathcal D$ for the rational
polarized reference space and a neat arithmetic group $\Gamma'$ of
lattice isometries containing the monodromy. Let
\[
                    \varphi:U_N\longrightarrow\Gamma'\backslash\mathcal D
\]
be the full period map. Assume that the normalization $B$ of the
closure of its actual image is an algebraic variety, that
the natural rational map $T\dashrightarrow B$ is birational, and
that $\varphi$ is generically immersive. In particular,
\[
                         \C(T)=\C(B).
\]
These assumptions describe the actual period-image model. The
algebraicity of $B$ is the period-image theorem of
\cite[Theorem~1.1]{BakkerBrunebarbeTsimerman}; here its construction
is an explicit input to the rank-loss statement.

\begin{theorem}[Boundary and ramification rank loss]\label{thm:rankloss}
For the data just specified, put $r=\nu(J)$.
Then $r>0$. If $E$ is either a component of $D_N$, or a component
of $D_T$ fixed pointwise by a nonidentity element of
$G_f$, then
\[
                           \nu(J|_E)<r.
\]
In particular, after placing the divisorial fixed loci in $D_T$
on an equivariant smooth model, the conclusion applies
to every divisorial ramification component of $T\to T/G_f$.
\end{theorem}

\begin{proof}
Proposition~\ref{altord:h:numerical-rank} identifies $r$ with the
generic rank of the highest-line map. We first show that this
rank is positive, then consider the nilpotent and finite cases.

If $r=0$, the highest line is constant on the connected universal
cover of $U_N$. Monodromy preserves it. Its connected algebraic
monodromy group $G$ is semisimple and hence acts trivially on that
line. Lemma~\ref{altord:p:normality}, applied with the normal
subgroup equal to $G$, gives $G=1$. After a finite \'etale cover the
monodromy is trivial, and the fixed-part theorem makes the entire
variation constant. Its full period map then has rank zero,
contrary to generic immersivity on the positive-dimensional variety
$T$. Thus $r>0$.

\smallskip\noindent
\emph{A component with nonzero monodromy logarithm.}
Let $E\subset D_N$ and write $N\ne0$ for its rational monodromy
logarithm. In coordinates $(u,v)$ near a general point of $E$,
with $E=(u=0)$, the lifted highest-line map has the form
\[
                [e^{zN}a(u,v)],\qquad z=\frac{\log u}{2\pi i}.
\]
Proposition~\ref{altord:h:boundary-line} provides an integer
$\ell\geq0$ with
\[
 N^{\ell+1}a(0,v)=0,\qquad N^\ell a(0,v)\ne0,
 \qquad
 \nu(J|_E)\leq\rk_{\rm gen}d\bigl(v\mapsto[N^\ell a(0,v)]\bigr).
\]
Lemma~\ref{altord:h:nilpotent-limit} gives its normal
tangential convergence on a sector from the interior line maps.
All interior $(r+1)$-minors vanish. Cauchy convergence of tangential
derivatives therefore bounds the rank of the limiting map by $r$.

Suppose that $\nu(J|_E)\geq r$. The limiting map has rank exactly
$r$ at a suitable general point, and its image lies in
$\mathbb P(\ker N)$. Lemma~\ref{altord:h:rank-principle} forces all
nearby interior lines to lie in this same linear subspace. The
condition that $N$ annihilates the highest line is holomorphic, so
it continues over the connected universal cover of $U_N$.

The one-parameter group $\exp(tN)$ belongs to $G$: it is the
connected Zariski closure of the cyclic unipotent local monodromy.
Every monodromy conjugate of this group fixes the highest line
pointwise. Let $G_N$ be the algebraic subgroup they generate.
It is connected, defined over $\Q$, and normalized by monodromy;
Zariski density makes it normal in $G$. Its invariant space contains
the highest line at a Hodge-generic point. By
Lemma~\ref{altord:p:normality}, $G_N=1$, contradicting $N\ne0$.
This proves the required inequality for components of $D_N$.

We have controlled all divisors at which the pure variation fails
to extend. The remaining obstruction is a divisor of finite
ramification where the variation does extend. Its inertia element
can act on the rational fibre even when it fixes every point of
the divisor, so we must use the full-period conclusion of
Lemma~\ref{altord:p:inertia}.

\smallskip\noindent
\emph{A divisor fixed by finite inertia.}
Let $b\in G_f$ be nonidentity and fix $E$ pointwise. If $E$ is
contained in $D_N$, the preceding case applies. Otherwise a general
point $e\in E$ lies in $U_N$. Choose a lift $h\in H$ of $b$.
The retained $H$-action supplies a flat Hodge isomorphism
\[
                  \Phi_y:V_y\longrightarrow V_{b(y)}
\]
on $U_N$. Use $V_e$ as a flat reference space on a small simply
connected $b$-invariant neighborhood, and set $a=\Phi_e$.
If $x$ is the full lifted period map, flatness gives
\[
                             x(by)=a\,x(y).
\]
The operator $a$ is a rational Hodge isometry at $e$. Because $H$
acts on the variation and $b(e)=e$, the group-action law gives
$a^{\operatorname{ord}(h)}=1$. Since $b$ fixes $E$, every highest line
on the nearby part of $E$ is an eigenline of $a$. Its eigenvalue
is constant on a connected small patch, so these lines lie in one
eigenspace $\ker(a-\lambda)$.

Proposition~\ref{altord:h:boundary-line} with $N=0$ bounds
$\nu(J|_E)$ by the rank of this extending line map on $E$.

Suppose again that $\nu(J|_E)\geq r$. The rank on $E$ is at most
$r$, since all $(r+1)$-minors of the ordinary extending line map
vanish by continuation from the interior. It is therefore $r$.
Apply Lemma~\ref{altord:h:rank-principle} in a transverse coordinate
with $N=0$. A nonempty open set of lifted highest lines lies in
$\ker(a-\lambda)$. Lemma~\ref{altord:p:inertia}, applied on the
smooth algebraic variety $U_N$ with fixed point $e$, now shows
that $a$ preserves every full period. Consequently $x(by)=x(y)$.

Projecting to $\Gamma'\backslash\mathcal D$ gives
$\varphi(by)=\varphi(y)$ on a nonempty open set. Since $B$ is the
normalization of the actual image, its normalization map is
birational and hence injective on a dense open subset of that
image. It follows that $b$ acts trivially on $\C(B)$. The chosen
model has $\C(T)=\C(B)$, so $b$ acts trivially on $T$, contrary
to its nonidentity in $G_f$. This proves the strict inequality.

Finally, in a finite quotient in characteristic zero, ramification
at a divisorial valuation is detected by a nontrivial inertia
group, which acts identically on that divisor's function field.
Its elements therefore fix the divisor pointwise. The last
assertion follows from the case just proved.
\end{proof}

\begin{remark}
The conclusion concerns the numerical dimension of the highest line,
while generic immersivity is a hypothesis on the full period map.
Neither the proof nor the conclusion identifies these ranks. The
function-field equality with the actual period image is used only
at the last contradiction in the finite case; it cannot be replaced
there by an arbitrary finite extension of that function field.
\end{remark}
